\scanpage{1}
% not manuscript: page de titre imprimée de l'Université de Montpellier / IMAG.

\scanpage{2}
% not manuscript: couverture portant « Pages 1-292 + table des matières ».

\scanpage{3}
\margin{I}

The ``long march'' through Galois theory...

1. Multi-Galois topoi \hfill 1

2. Applications to coverings of topoi \hfill 4

3. Pro-multi-Galois variants \hfill 6

4. Complements remorse

\figure{[Brace on the right joining entries 3 and 4.]}

\uuline{5.} Introduction of the arithmetic context
Fundamental ``anabelian conjecture'' \hfill \(7'\)

6. Local analysis of \((X,S)\) at an \(s\in S\). \hfill \(10'\)

7. ``Messy'' reformulation of the conjecture
(the necessary purgatory...) \hfill 12

\deleted{\unsure{6.}}
8. Taxonomic reflection (distinction \unsure{of the|of}
cases, \unsure{where} purgatory gets fitted out a little...) \hfill \(25'\)

\deleted{\unsure{7.}}
9. Tangential structure at the \(s\in S\)
(Sections of extensions ``of the second type'' \hfill 34

10. \unsure{Adjustment} of the hypotheses, (remorse) \hfill 38

11. Conditions on the systems of groupoids obtained from
geometric situations \hfill 39

12. The topological analogy. (Where we also convince ourselves that the groupoidal
mess can be expressed completely, in the anabelian cases,
by the outer groups with loops... \hfill \unsure{41|43}

13. Return to the arithmetic case: ``Galois'' formulations \hfill 53

13\textsuperscript{bis}. Return to the notion of a group with loops \hfill 56

14. Cohomological digression (on ``plugging holes'') \hfill 58

\uuline{15.} Return to the topological case: critical orbits
\unsure{around} splittings of extensions
Application to finite subgroups of
\(\mathrm{Autext}_{\mathrm{lac}}(\pi)\) (discrete case)
(cf. also \(\S 18\)) \hfill \(67'\)

16. Plugging and drilling holes: general topological preliminaries
\hfill \(79'\)

17. Complements to \(\S 15\), on \unsure{operations|groups} of finite groups
on surfaces \hfill 89

\scanpage{4}
\uuline{18.} Drilling holes. Application to finite
subgroups of \(\mathrm{Autext}_{\mathrm{lac}}(\pi)\)
(discrete case) \hfill 91

\uuline{19.} Teichmüller tower \hfill 102

20. Digression: 2-isotopic description
of the category of topological \unsure{2-spheres} \hfill \(116'\)

21. Teichmüller spaces \hfill 125

\margin{missing \(\S 22\)!}

23. Return to surfaces with (finite) groups of operators
(``putting into equations'' the problem) \hfill 133

24. Attempt to determine the \unsure{\(A^{\circ\Gamma}\)}. Link with the relations
\(\Pi_{g,(\nu,\nu+n-1)}^{\Gamma}\) \(=(1)\),
program of \deleted{\illegible{1}}\inserted{work} \hfill 142

25. ``\uline{Special}'' Teichmüller groups \hfill 148

25\textsuperscript{bis}. ``Case of the two groups'' of operators. Return to the notation.
\hfill \(153'\)

\uuline{25.} Profinite Teichmüller groups. \hfill 161

Discretifications and \unsure{prediscretifications|prodiscretifications}. Link with modular
Teichmüller topos. Hasty conjecture. \hfill 167

27. Change of type \((g,\nu)\): a) plugging holes
(and \unsure{remarkable diagrams|remarkable diagram} page 174) \hfill 172

28. change of type \((g,\nu)\): b) passage to a finite covering
(the hasty conjecture creaks...) \hfill 177

29. Critique of the preceding approach. (We readjust
the notions and the conjectures...) \hfill 186

30. Properties of the \(N_{g,\nu},\Gamma_{g,\nu}\): a)
\unsure{Properties related} to finite Teichmüller subgroups.
\hfill 192

31. Digression on liftings of an outer action of a finite
group on a profinite group with loops \hfill 198

\uuline{32.} Return to the arithmetic aspects of plugging holes:
relations between \(\Gamma_{g,\nu}\) and \(\Gamma_{g,\nu-1}\).
(Where we reconstruct, without saying so, \uline{the whole étale topos} of a
\inserted{complete} algebraic curve, from the \(\pi_1\) of an
anabelian open...) \hfill 200

\uuline{33.} Topological digression: anti-involutions of compact
oriented surfaces (leaving aside at first the case ``with holes'')
\hfill 221

\uuline{\unsure{35.}} The injectivity of
\(\Gamma_{\mathbb Q}\longrightarrow
\mathrm{Autext}_{\mathrm{lac}}(\widehat{\Pi}_{0,3})\)
\hfill 249

\unsure{36.} The isomorphism
\(\Gamma_{\mathbb Q}\simeq\Gamma_{1,1}\),
and the injectivity of
\(\Gamma_{\mathbb Q}\longrightarrow
\mathrm{Autext}(\widehat{\unsure{\Gamma|T}}_{1,1}^{+})
=\mathrm{Autext}(\mathrm{SL}(2,\mathbb Z)^{\wedge})\)
\hfill 257

\unsure{37.} Moduli of elliptic curves via Legendre, or
\unsure{\(M_{1,1}[2]'\simeq U_{0,3}\ne M_{0,4}^{!}\)} \hfill 272

33. Relations between \(N_{g,\nu},\Gamma_{g,\nu}\) for variable \(g\)
(study of finite coverings) \hfill 239

34. Heuristic profinite description of the category of algebraic curves
defined over finite \deleted{\unsure{extensions}}
subextensions of \(\mathbb Q_0/\mathbb Q\)
\hfill 245

\figure{[To the left of the last two entries, brace extended
upward by an arrow pointing to the entry
``The injectivity of \(\Gamma_{\mathbb Q}\)''.]}

\scanpage{5}
\margin{\unsure{814|714}\par\unsure{752|252}\par 832}

\scanpage{6}
\mspage{1}
\margin{\unsure{W|VI}}

1. Multi-Galois topoi.

\uline{Proposition} 1.1. Let \(\mathcal E\) be a category.
Equivalent conditions:

a) \(\mathcal E\) is a topos, and every object of \(\mathcal E\)
is local\textsuperscript{ly} constant.

b) \(\mathcal E\) is equivalent to a category \(\widehat C\),
where \(C\) is a groupoid.

\margin{N.B. We shall see (below) \unsure{that we can choose}
\unsure{\(C\)} \uline{canonically}}

b') \(\exists\) a family \((G_i)_{i\in I}\) of groups
and an equivalence of categories
\[
\mathcal E\approxeq\prod\mathrm{Ens}(G_i)
\]

c) Ad-hoc exactness conditions, of the type of those given in SGA 1...

Proof. \(b\Longleftrightarrow b'\Longrightarrow a\) immediate.
For \(a)\Longrightarrow b)\), I am less sure, perhaps one has to
assume that \(\mathcal E\) is locally connected, and that it has sufficiently
many fiber functors i.e. sufficiently many points.

\uline{Definition 1.2.} If the equivalent conditions b) b') above
are satisfied, one says that \(C\)
is \unsure{a|a} multi-Galois \deleted{category} \inserted{topos}
(or a multi-Galois \uline{category}).

\margin{\uline{Proposition 1.3}
a) If \(\mathcal E\) is multi-Galois, every induced \unsure{topos}
\(\mathcal E_{/S}\) too.
\par
b) every sum of multi-Galois topoi (i.e. \deleted{\illegible{1}}
product of multi-Galois categories) \unsure{is so|is} too.}

\scanpage{7}
\uline{Proposition 1.4} \unsure{Let|Let} \(\mathcal E\) be a topos,
\(C\) the category of points of \(\mathcal E\) (opposite to the
category of fiber functors on \(\mathcal E\)). The canonical
functor
\[
\mathcal E\times C^o\longrightarrow(\mathrm{Ens})
\]
induces a canonical functor
\[
\mathcal E\longrightarrow\mathrm{Hom}(C^o,(\mathrm{Ens}))
\overset{\mathrm{déf}}{=}\widehat C
\]

This being so: (\(\mathcal E\) being multi-Galois)

a) \(C\) is a groupoid (called the ``fundamental groupoid'' of the
multi-Galois topos \(\mathcal E\), and often denoted
\(\Pi_1(\mathcal E)\)).

b) \(\mathcal E\longrightarrow\widehat C\) is an equivalence
of categories.

An object \inserted{\(S\)} of a topos \inserted{\(\mathcal E\)}
is said to be 0-connected if it is \(\ne\varnothing_{\mathcal E}\)
(\unsure{it|i.e.} is not an initial object) and if it is connected
(i.e. \(S\simeq S'\amalg S''\) implies
\(S'\) or \(S''\simeq\varnothing_{\mathcal E}\)) --- this also means
that the induced topos \(\mathcal E_{/S}\) is 0-connected
i.e. is not the initial topos (``empty topos'', i.e.
\inserted{equivalent to the} terminal category) and that it is connected,
i.e. ... One says that \(S\) is 1-connected

\scanpage{8}
\mspage{2}

(or ``simply connected'') if it is 0-connected, and if every
\deleted{\unsure{covering of \(S\) is}}
\inserted{object \(S'\) of \(\mathcal E_{/S}\)
\deleted{\illegible{1}}}
locally constant \unsure{is} is \unsure{\deleted{constant}|constant}
--- which depends \inserted{again} only on the induced topos
\(\mathcal E_{/S}\), which will then be said to be 1-connected.

\uline{Proposition} 1.5. Let \(\mathcal E\) be a multi-Galois topos,
and let \(S\) be an object of \(\mathcal E\). Equivalent conditions

a) \(S\) is 1-connected

b) \(S\) \unsure{(is)|is} 0-connected and projective

c) the covariant functor represented by \(S\)
\[
T\longmapsto\mathrm{Hom}_{\mathcal E}(S,T)
\]
is a fiber functor, or again
as it is already left exact) it commutes with
arbitrary \(\underset{\longrightarrow}{\lim}\)

\inserted{N.B.} (it suffices that it commute with sums and with passage
to the quotient...)

d) \(\mathcal E_{/S}\) is equivalent to the point topos

e) (If \(\mathcal E=\widehat C\), \(C\) a groupoid) the functor
\(S\) on \(C\) is \inserted{representable}

\scanpage{9}
\uline{Definition 1.6} One then says, sometimes, that \(S\) (or better, the topos
\(\mathcal E_{/S}\)) is a \uline{universal} \uline{covering} of the
\inserted{multi-Galois} topos \(\mathcal E\).

\uline{Proposition 1.7}
\deleted{\unsure{Taking} \illegible{1} object}
\inserted{Let \(\mathcal E_0\) be} the full subcategory of
\(\mathcal E\) formed by the 1-connected objects
(of the multi-Galois category \(\mathcal E\)),
\deleted{Then \unsure{the functor}} \inserted{and \(\Pi_1(\mathcal E)\)
the fundamental groupoid} of \(\mathcal E\).
We have \inserted{\unsure{by} 1.5} a canonical functor
\[
\mathcal E_0\longrightarrow
\Pi_1(C)
\]

(associating: \(\forall\) \(S\in\mathrm{Ob}\,\mathcal E_0\)
the fiber functor it represents, or rather the corresponding ``point''
of \(C\)), which is
(not only fully faithful, but even) an equivalence of
categories: every fiber functor on \(\mathcal E\) is representable
(by an essentially unique (1-connected) object, as is only right...).

\scanpage{10}
\mspage{3}
\deleted{\unsure{Definition 1.8} \illegible{1} \unsure{The object}}

\uline{Corollary 1.8.} Let \(P\) be a ``point'' of \(\mathcal E\)
\inserted{(associated with a fiber functor \(F_P\))}.
There exists a 1-connected object \(S\) of \(\mathcal E\) and a lifting
\(P\xrightarrow{\alpha}\mathcal E_{/S}\)
(i.e. an \(\alpha\in F(S)\)),

\diagram{[Under \(P\xrightarrow{\alpha}\mathcal E_{/S}\),
\(\mathcal E\) is placed at bottom center. Two dotted arrows
go respectively from \(P\) and from \(\mathcal E_{/S}\) to
\(\mathcal E\). Between them, a rightward arrow with
\(\sim\) above it, \unsure{double-lined|single-lined}.
Two side brackets enclose this lower part of the diagram.]}

and this determines \((S,\alpha)\) up to unique isomorphism.

In fact, \(S\) is the unique \unsure{object|\deleted{object}}
of \(\mathcal E\) that represents \(F_P\)...

\uline{Definition 1.9.} One says that \(S\)
\inserted{(or \(\mathcal E_{/S}\))} is the
\uline{universal covering pointed over \(P\)} determined
by the point \(P\).

\uline{Scholium 1.10} Giving oneself a ``point'' of the \unsure{:}
\inserted{multi-Galois} topos \(\mathcal E\), or giving oneself a
universal covering, amounts essentially to the
same: each determines the other...
